\section{Discussion}\label{sec:discussion}

\subsection{Theory, predictions and results}\label{sec:mapping}

Table~\ref{tab:mapping} sets each hypothesis beside the prediction it made in our variables, the test, the result and our reading of it. The principal--principal frame survives the four tests in the form we gave it. The filing is informative about the controller's intention (H1), and informative in a way that does not depend on the models knowing the companies. The controller's stake is associated with what the filing signals and, through the signal, with what happens (H2). The block predicts the signal, and once the signal is in the regression the block adds nothing to the prediction of early disposal; that is what we would expect if the stake works by shaping the intention that the filing discloses rather than through some channel the filing does not capture. The market does not read the filing as a transfer (H3). The announcement reaction rises with the signal, the subsequent return falls with it, the fall is realised in the parents that dispose early and in those with small blocks, and the signal has no explanatory power for the return once realised early disposal is known. The reform changed the filings and their reception and, so far, neither the disposal signal nor the rate of early disposal (H4); captive spin-offs also became rarer, a change we record without attributing it.

Two of the readings deserve the qualification the table gives them. The association between the controlling block and the realised disposal is clear in the tercile sort (\LowBlockDisp{} against \HighBlockDisp{} percent) and imprecise in the regression with controls (\BlockOnDisposed{}, standard error \BlockOnDisposedSE). What is precise is the block's association with the signal and with the return. The reform result is a before-and-after comparison with \NPostFilings{} post-reform filings and a short outcome record. It establishes that the disposal signal and the early-disposal rate did not fall over the horizon we observe, not that the appraisal right and the listing review will not change behaviour once they have been exercised against a listing.

\begin{sidewaystable}[htbp]
\centering
\caption{Theory, prediction, test and result}
\label{tab:mapping}
{\footnotesize\setlength{\tabcolsep}{3pt}% generated by code/make_tables_cgir.py on 2026-10-26 — source: all results files (see rows); readings recorded by the author 2026-10-25
\begin{tabular}{>{\raggedright\arraybackslash}p{0.9cm}>{\raggedright\arraybackslash}p{3.4cm}>{\raggedright\arraybackslash}p{3.6cm}>{\raggedright\arraybackslash}p{3.3cm}>{\raggedright\arraybackslash}p{6.2cm}>{\raggedright\arraybackslash}p{2.4cm}}
\toprule
 & Theoretical statement & Prediction in our variables & Test & Result & Reading \\
\midrule
H1 & Disclosure is informative about the controller's intention & The filing text predicts disposal better than the parent's financials & Table~\ref{tab:gate1}: AUC two-model vs financial baseline, pooled; anonymized condition & AUC 0.761 vs 0.587 (difference 0.174, 95\% CI [0.117, 0.235]); anonymized 0.691 (2001--16) and 0.803 (2017+) & Supported; not memorization \\ \addlinespace[2pt]
H2 & Smaller controlling blocks make disposal more attractive to the controller & Lower block $\Rightarrow$ stronger disposal signal and more frequent (early) disposal & Tables~\ref{tab:block_terciles}, \ref{tab:block_regressions}: score and disposal on the block & Score on block $-$0.113 (HC1 $p$ 0.013); disposal rate 74\% vs 62\% (low vs high tercile); early disposal on block 0.113 ($p$ 0.395) & Signal: supported; realized disposal: direction only \\ \addlinespace[2pt]
H3a & Minority investors read the filing (rival: they do not) & If read: more negative announcement CAR for high-signal filings and no later drift & Table~\ref{tab:returns_by_score}: announcement CAR and one-year BHAR on the score & Announcement slope 0.037 ($p$ 0.046) --- positive; one-year Spearman $-$0.136 ($p$ 0.001); within-year Q5$-$Q1 $-$11.7\% & Rival supported: reaction positive, drift negative \\ \addlinespace[2pt]
H3b & The loss arrives with the disposal & Early disposers underperform retainers; late disposers do not within the year & Tables~\ref{tab:returns_by_outcome}, \ref{tab:realized_controls} & Early $-$ never $-$15.6\% (Mann--Whitney $p$ 0.001); late $-$ never $-$5.2\% ($p$ 0.626); score adds nothing once early disposal is in (Table~\ref{tab:realized_controls}) & Supported \\ \addlinespace[2pt]
H3c & Minority holders of low-block parents bear the loss & Lower block $\Rightarrow$ lower one-year return & Tables~\ref{tab:block_terciles}, \ref{tab:block_regressions} & Median BHAR $-$26.9\% (low) vs $-$9.1\% (high), Mann--Whitney $p$ $<$0.001; rank on block 0.142 ($p$ 0.062) & Supported (association) \\ \addlinespace[2pt]
H4 & Mandated disclosure changes what is said, not what is done & Protection statements and announcement reactions change; the disposal signal and early disposal do not & Table~\ref{tab:reform} & Protection statements 0.5\% $\to$ 95.5\%; announcement CAR +0.06\% $\to$ +5.29\%; score 0.528 $\to$ 0.510; early disposal 40.2\% $\to$ 36.8\% & Consistent (before/after, small post sample) \\ \addlinespace[2pt]
\bottomrule
\end{tabular}
}
\tabnote{Each row states the theoretical claim behind a hypothesis of Section~\ref{sec:theory}, the prediction it makes in the paper's variables, the table that tests it, the result read from that table, and our reading. ``Supported'' means the association predicted is present with the predicted sign and conventional significance; ``consistent'' means the pattern is present but the design (a before-and-after comparison without a control group, or a small subsample) does not permit a stronger word. None of the rows is a causal estimate.}
\end{sidewaystable}

Table~\ref{tab:mirror} places our results beside those of \citeA{kang2025}. We did not use their sample, our period extends theirs by five years at the end, and our measures differ from theirs in ways the table states. The two papers agree on the facts they both measure: a small positive announcement return, a negative return over the following year that is concentrated in the disposed subsidiaries and in the parents with small controlling stakes, and a disposal rate near seventy percent. They differ on the announcement return's relation to ownership, which is positive and significant in their sample and positive and insignificant in ours. And they differ in what they test. \citeA{kang2025} establish the loss and locate it. We establish that the loss was forecastable from the controller's own disclosure at the time the minority still held the asset, and that the market's reaction to that disclosure had the wrong sign.

\begin{table}[htbp]
\centering
\caption{The tests of \protect\citeA{kang2025} and their counterparts in this paper}
\label{tab:mirror}
{\footnotesize\setlength{\tabcolsep}{3pt}% generated by code/make_tables_cgir.py on 2026-10-28 — source: literature/anchor_numbers.json (Kang & Lee 2025, Table 1 Panel B, Tables 2-3, footnotes 14 and 18, Section 4) and our results files
\begin{tabular}{>{\raggedright\arraybackslash}p{5.2cm}>{\raggedright\arraybackslash}p{4.6cm}>{\raggedright\arraybackslash}p{5.4cm}}
\toprule
Test & \protect\citeA{kang2025}, 525 captive spin-offs 2001--2021 & This paper \\
\midrule
Announcement CAR [$-$1, +1] & 0.853\% ($t$ = 2.31) & +0.58\% ($t$ = 1.3), $n$ = 583 \\
One-year abnormal return after the announcement & CAR[+2, +252] $-$7.857\%; BHAR[+2, +252] $-$10.995\% (matched portfolio) & BHAR[+2, +250] mean $-$6.0\%, median $-$15.9\% (market-adjusted), $n$ = 578 \\
Share of subsidiaries eventually disposed of & 69.3\% (364 of 525) & 69.3\% of 579 clean-label events; listed share 3.3\% of 581 clean labels \\
Long-run return, disposed vs retained & disposed $-$16.7\% to $-$10.45\%; retained slightly positive, insignificant & disposed $-$8.9\% vs retained +1.4\% (Mann--Whitney $p$ 0.010); early disposers $-$14.2\% \\
Controlling ownership and long-run return & below-median ownership $-$19.38\% to $-$12.18\%; above-median $-$2.58\% to $-$1.01\%; coefficient 0.511--0.769 & low tercile median $-$26.9\% vs high $-$9.1\%; rank on block 0.142 ($p$ 0.062) \\
Controlling ownership and announcement CAR & coefficient 0.052** & 0.030 (HC1 $p$ 0.199) \\
Ownership and the retention decision & ownership positively associated with retention (unreported) & disposal rate 74\% (low block) vs 62\% (high); disposal on block $-$0.092 ($p$ 0.474) \\
Can investors discern the intention at the announcement? & stated as unlikely; not tested & tested: the filing predicts disposal (Table~\ref{tab:gate1}); the announcement reaction does not price it (Table~\ref{tab:returns_by_score}) \\
\bottomrule
\end{tabular}
}
\tabnote{Left column: results as reported by \citeA{kang2025} for \KLN{} captive spin-offs announced by listed parents between 2001 and 2021 (their Tables 1--3, Section 4 and footnotes 14 and 18); their long-run abnormal returns use a matched portfolio and their ownership measure is the controlling family's and inside directors' stake. Right column: this paper's counterparts, with market-adjusted buy-and-hold returns and the controlling block defined in Table~\ref{tab:defs}. The two samples were built independently from the regulator's filings.}
\end{table}

\subsection{Theoretical implications}\label{sec:theoretical}

The principal--principal literature has treated the controlling owner's intention as unobservable and has inferred it from outcomes: a merger that moves value toward the family's firms, an offering priced below market, a loan to a weak affiliate \cite{bertrand2002, bae2002, baek2006, buchuk2014}. The inference is sound but late, because the outcome is the expropriation. Our results say that for one class of transactions the intention was observable before the outcome, in the document the controller was legally required to produce to obtain the minority's vote. It was observable not because the controller confessed but because a disposal-oriented division has to be described differently from a retention-oriented one if the description is to serve the controller's later purposes. This moves the question the literature asks about disclosure. The question has been whether more disclosure deters expropriation \cite{he2025}. The prior question is whether the disclosure already produced contains the expropriation's forecast, and whether anyone reads it. In our setting the disclosure contained the forecast, and the announcement reaction moved in the direction opposite to it (a slope of \AnnCarSlope{} on the signal, Table~\ref{tab:returns_by_score}), so the question is not whether the price missed a small effect but why it priced the signal with the wrong sign.

The result also qualifies the efficient-markets presumption that attaches to event studies of divestitures. The carve-out literature in the United States reads a positive announcement return as the market's assessment of the restructuring's value \cite{schipper1986, nanda1991, slovin1995, hulburt2002}. Read that way, the Korean announcement returns in Table~\ref{tab:returns_by_score} would say that the market welcomed the divisions most likely to end in a disposal. The following year says otherwise. In a market in which the controlling block is large and the disposal is at the controller's discretion, the announcement return measures what minority investors understood, and what they understood, on this evidence, was not the disposal the filing's details implied. The forecast is extractable by a reader with no stake in the company and no knowledge of its later history, and it was not extracted by the investors who did have a stake. That is a fact about attention and expertise in minority ownership, of the kind \citeA{anand2026} document when retail investors' governance proposals respond to the readability of the 10-K, and one the principal--principal frame has not had to confront, because it has not previously had a measure of the intention separate from the outcome.

The controlling block's role is the one \citeA{young2008} and \citeA{aguilera2025} describe, but the chain we trace is specific. The block does not move the announcement reaction and does not, once the signal is in the regression, move the probability of early disposal; it moves the signal, and the signal predicts the disposal. The controller's incentive, in other words, shows in what the controller writes before it shows in what the controller does. A theory of controlling-owner behaviour that wants to predict expropriation from ownership structure can, for transactions that are announced before they are completed, add the announcement to its data.

\subsection{Practitioner implications}\label{sec:practitioner}

For the minority investor the implication is direct and uncomfortable. The information needed to avoid the typical loss was in the filing, and extracting it required reading the filing for what it implied about the new company's future ownership rather than for its stated purpose. A reader who ranked filings by the signal we measure would have placed in the top quintile a group of parents whose median one-year market-adjusted return was \QFiveMedian{} percent, against \QOneMedian{} percent in the bottom quintile. A reader who took the announcement reaction as the market's verdict would have been reassured by the top quintile's reception and wrong a year later. The signal is strongest for the disposals that follow within a year, which is also where the loss is. It is a reading of the filing's description of the new company's future ownership rather than of its stated purpose, which the post-reform filings show can be made uniform by a form.

For boards and controlling owners the finding is that the filing is read, if not by the market then by anyone who later wishes to show what the controller knew. From October 2022 the filing must state the division's purpose and the protection offered to shareholders, and the exchange will read it when the subsidiary applies to list. A controller who intends to dispose of the subsidiary and writes a retention-oriented filing creates a record that the listing review can hold against the parent. A controller who writes candidly has, on our evidence, nothing to fear from the announcement reaction, which has so far received the filings that read as disposals more favourably than the others. Post-reform filings state protection measures and listing plans at rates (\ProtectPost{} and \ListingPost{} percent) that leave the statements with little power to distinguish one filing from another, and the announcement reaction to them has been positive. We would not expect the statements to stay uninformative once the listing review has been exercised against a parent whose statements were not honoured.

Institutional investors and proxy advisers vote on the division and not on the disposal. The division vote is therefore the only point at which the disposal can be contested, and the filing is the only document on which to contest it. Our measure is a reading of that document that a fund's own analysts could replicate. The dimension it captures, the implied future ownership of the new company, is the one on which engagement before the vote would have to turn.

\subsection{Policy implications}\label{sec:policy}

The disclosure measure of the 2022 reform rests on the premise that general shareholders lacked the information to decide. In the Commission's own description, the mandated statements are to give them sufficient information for the vote, the appraisal right and their investment decisions \cite{fsc2022b}. Our evidence says that the information that mattered was already in the filing and was not priced, so the disclosure measure addressed a problem that was not the binding one. The two measures that alter the controller's incentives are the appraisal right, which makes the division costly if minority holders dissent in numbers, and the listing review, which can refuse the disposal's most valuable form. Those are the measures that act on the mechanism we document, and they are the ones whose effects cannot yet be seen in the data. Disclosure rules can move shareholder behaviour when they change what shareholders can see about each other (the Japanese requirement to disclose institutions' voting records raised dissent in director elections; \citeNP{uchida2025}), but the 2022 statements are addressed to shareholders who, on our evidence, were not reading the filing to begin with. The practical reading for the Financial Services Commission and the Korea Exchange is that the content of the mandated statements should be audited against subsequent conduct. A protection statement in a filing whose subsidiary applies to list within the five years the review covers should be compared with what the parent did, and the comparison published. A disclosure requirement that produces statements in \ProtectPost{} percent of filings and no change in the behaviour disclosed is a form, not a safeguard.

The appraisal right must be exercised at the division, at a price based on the market before the division was pursued. At that moment the market's reaction is positive, so a shareholder who takes the price as the verdict has no reason to dissent, and the loss we document arrives over the following year as the disposal materialises. The right as designed protects the shareholder who has read the filing as our models read it, not the one who has read the market. A right exercisable at the disposal rather than at the division, or a requirement that the parent's shareholders vote on a disposal of the subsidiary within some years of a captive spin-off, would act at the step where the loss occurs. \citeA{kang2025} reach the same recommendation from the outcome evidence; the forecastability we add implies that such a vote would be informed. More generally, regulators of concentrated-ownership markets face other reorganisations with a deferred and discretionary second step: a transfer of a business to an affiliate ahead of its sale, a holding-company conversion. For these the finding is that the first step's disclosure can be read for the second step's probability, and that the minority's protection should be attached to the second step.
